% Figure S3 legend — Nature-style supplementary figure caption
% Companion to: scripts/supplementary_figures/figure_s3.py

\textbf{Figure S3 $|$ Concordance between scRNA-seq and bulk RNA-seq expression
profiles.}

\textbf{(A)}~Per-gene mean expression in scRNA-seq versus bulk RNA-seq across three
matched conditions, plotted on log-log axes.
Each point represents one gene; only genes expressed above background in both
datasets are shown ($n$ indicated per panel).
Pearson $r$ is computed on $\log_{10}$-transformed values; Spearman $\rho$ on raw
values.
Exponential-phase cells: $r = 0.63$, $\rho = 0.61$ ($n = 1{,}239$ genes).
Dis-Arrest cells: $r = 0.75$, $\rho = 0.75$ ($n = 1{,}684$ genes).
Reg-Arrest cells: $r = 0.34$, $\rho = 0.23$ ($n = 916$ genes).
The weaker concordance for Reg-Arrest reflects differences in the bulk antibiotic
condition used as a proxy; the LFC comparison in panels~B--C, which relies on the
same pseudobulk contrast, yields stronger agreement (see below).

\textbf{(B,~C)}~Per-gene log$_2$ fold-change (LFC) compared between the sc pseudobulk
DESeq2 analysis and the bulk RNA-seq DESeq2 analysis, for Dis-Arrest~\textbf{(B)}
and Reg-Arrest~\textbf{(C)} versus the exponential-phase control.
The pseudobulk analysis aggregates per-condition single-cell counts into
replicate-level profiles and applies the same DESeq2 pipeline used for bulk
data~(Methods).
Correlation coefficients: Dis-Arrest $r = 0.62$, $\rho = 0.61$;
Reg-Arrest $r = 0.33$, $\rho = 0.31$ ($n = 3{,}466$ genes in both panels).
The positive correlation in both conditions confirms that the sc pseudobulk
and bulk experiments capture the same genome-wide transcriptional response.

\textbf{(D)}~GO enrichment of down-regulated genes identified by the sc pseudobulk
DESeq2 analysis (top 500 genes by LFC magnitude, matched study sets).
Bars show $-\log_{10}(\text{FDR})$ for GO Biological Process terms that reach
significance (FDR~$<$~0.05) in both Dis-Arrest (red) and Reg-Arrest (blue).
The same functional categories suppressed in the bulk analysis (translation,
flagellar motility, chemotaxis, energy metabolism) are independently recovered
in the single-cell data, validating the biological signal detected across both
platforms.
Reg-Arrest displays stronger enrichment than Dis-Arrest for most terms,
consistent with the bulk GO results in Figure~4B.
